\subsection*{Morphisms}
\label{subsec:morphisms}

\begin{definitionbox}{Definition (Homomorphism)}
\begin{definition}[Homomorphism]
\label{def:algebraic-homomorphism}
Let $A$ and $B$ be algebraic structures of the same type.
A \textbf{homomorphism} $f:A\to B$ is a function that preserves the operations
of the structure.
\end{definition}
\end{definitionbox}

\begin{remark*}[Standard quantified statement]
\[
\operatorname{Homomorphism}(f,A,B)
\Longleftrightarrow
f:A\to B
\text{ and } f \text{ preserves the named operations of }A\text{ and }B.
\]
\end{remark*}

\begin{remark*}[Definition predicate reading]
$f$ is a homomorphism exactly when it is a structure-preserving function.
\end{remark*}

\begin{remark*}[Negated quantified statement]
\[
\neg\operatorname{Homomorphism}(f,A,B)
\]
when $f$ is not a function $A\to B$ or fails to preserve at least one named
operation.
\end{remark*}

\begin{remark*}[Interpretation]
Homomorphisms are the maps that keep algebraic operations visible through a
function.
\end{remark*}

\begin{dependencies}
\begin{itemize}
  \item \hyperref[def:alg-function]{Function}
  \item \hyperref[def:binary-operation]{Binary Operation}
\end{itemize}
\end{dependencies}

\begin{definition}[Monomorphism]
\label{def:monomorphism}
A \textbf{monomorphism} is an injective homomorphism.
\end{definition}

\begin{remark*}[Standard quantified statement]
\[
\operatorname{Monomorphism}(f,A,B)
\Longleftrightarrow
\operatorname{Homomorphism}(f,A,B)
\land
\operatorname{Injective}(f,A,B).
\]
\end{remark*}

\begin{remark*}[Definition predicate reading]
A monomorphism is a homomorphism that embeds without collisions.
\end{remark*}

\begin{remark*}[Negated quantified statement]
\[
\neg\operatorname{Monomorphism}(f,A,B)
\Longleftrightarrow
\neg\operatorname{Homomorphism}(f,A,B)
\lor
\neg\operatorname{Injective}(f,A,B).
\]
\end{remark*}

\begin{remark*}[Interpretation]
Monomorphisms preserve structure while keeping distinct inputs distinct.
\end{remark*}

\begin{dependencies}
\begin{itemize}
  \item \hyperref[def:algebraic-homomorphism]{Homomorphism}
  \item \hyperref[def:alg-injective]{Injective Function}
\end{itemize}
\end{dependencies}

\begin{definition}[Epimorphism]
\label{def:epimorphism}
An \textbf{epimorphism} is a surjective homomorphism.
\end{definition}

\begin{remark*}[Standard quantified statement]
\[
\operatorname{Epimorphism}(f,A,B)
\Longleftrightarrow
\operatorname{Homomorphism}(f,A,B)
\land
\operatorname{Surjective}(f,A,B).
\]
\end{remark*}

\begin{remark*}[Definition predicate reading]
An epimorphism is a homomorphism that reaches every element of the codomain.
\end{remark*}

\begin{remark*}[Negated quantified statement]
\[
\neg\operatorname{Epimorphism}(f,A,B)
\Longleftrightarrow
\neg\operatorname{Homomorphism}(f,A,B)
\lor
\neg\operatorname{Surjective}(f,A,B).
\]
\end{remark*}

\begin{remark*}[Interpretation]
Epimorphisms preserve structure while covering the whole target.
\end{remark*}

\begin{dependencies}
\begin{itemize}
  \item \hyperref[def:algebraic-homomorphism]{Homomorphism}
  \item \hyperref[def:alg-surjective]{Surjective Function}
\end{itemize}
\end{dependencies}

\begin{definition}[Isomorphism]
\label{def:isomorphism}
An \textbf{isomorphism} is a bijective homomorphism.
\end{definition}

\begin{remark*}[Standard quantified statement]
\[
\operatorname{Isomorphism}(f,A,B)
\Longleftrightarrow
\operatorname{Homomorphism}(f,A,B)
\land
\operatorname{Bijective}(f,A,B).
\]
\end{remark*}

\begin{remark*}[Definition predicate reading]
An isomorphism is a structure-preserving bijection.
\end{remark*}

\begin{remark*}[Negated quantified statement]
\[
\neg\operatorname{Isomorphism}(f,A,B)
\Longleftrightarrow
\neg\operatorname{Homomorphism}(f,A,B)
\lor
\neg\operatorname{Bijective}(f,A,B).
\]
\end{remark*}

\begin{remark*}[Interpretation]
Isomorphisms identify two structures as the same up to relabeling.
\end{remark*}

\begin{dependencies}
\begin{itemize}
  \item \hyperref[def:algebraic-homomorphism]{Homomorphism}
  \item \hyperref[def:alg-bijective]{Bijective Function}
\end{itemize}
\end{dependencies}

\begin{definition}[Endomorphism]
\label{def:endomorphism}
An \textbf{endomorphism} is a homomorphism from a structure to itself.
\end{definition}

\begin{remark*}[Standard quantified statement]
\[
\operatorname{Endomorphism}(f,A)
\Longleftrightarrow
\operatorname{Homomorphism}(f,A,A).
\]
\end{remark*}

\begin{remark*}[Definition predicate reading]
An endomorphism is a structure-preserving self-map.
\end{remark*}

\begin{remark*}[Negated quantified statement]
\[
\neg\operatorname{Endomorphism}(f,A)
\Longleftrightarrow
\neg\operatorname{Homomorphism}(f,A,A).
\]
\end{remark*}

\begin{remark*}[Interpretation]
Endomorphisms are internal symmetries or transformations, not necessarily
invertible ones.
\end{remark*}

\begin{dependencies}
\begin{itemize}
  \item \hyperref[def:algebraic-homomorphism]{Homomorphism}
\end{itemize}
\end{dependencies}

\begin{definition}[Automorphism]
\label{def:automorphism}
An \textbf{automorphism} is a bijective endomorphism.
\end{definition}

\begin{remark*}[Standard quantified statement]
\[
\operatorname{Automorphism}(f,A)
\Longleftrightarrow
\operatorname{Endomorphism}(f,A)
\land
\operatorname{Bijective}(f,A,A).
\]
\end{remark*}

\begin{remark*}[Definition predicate reading]
An automorphism is an invertible structure-preserving self-map.
\end{remark*}

\begin{remark*}[Negated quantified statement]
\[
\neg\operatorname{Automorphism}(f,A)
\Longleftrightarrow
\neg\operatorname{Endomorphism}(f,A)
\lor
\neg\operatorname{Bijective}(f,A,A).
\]
\end{remark*}

\begin{remark*}[Interpretation]
Automorphisms are the reversible symmetries of a structure.
\end{remark*}

\begin{dependencies}
\begin{itemize}
  \item \hyperref[def:endomorphism]{Endomorphism}
  \item \hyperref[def:alg-bijective]{Bijective Function}
\end{itemize}
\end{dependencies}
